\newpage
\section{Duração e Cronograma}

\subsection{Análise da Duração do Projeto}

O desenvolvimento do sistema \textbf{Tropical Mix} foi estruturado em quatro fases principais, abrangendo desde o planejamento e definição técnica até a entrega da aplicação final. O cronograma foi planejado para garantir um avanço progressivo, permitindo a validação de cada etapa antes do início da seguinte.  

\subsubsection{Fase 1 – Planejamento e Desenho da Aplicação (08/2025 – 09/2025)}

Esta fase abrangeu o levantamento de requisitos, definição das tecnologias e elaboração da arquitetura da aplicação. Também foram desenvolvidos os diagramas de modelagem, o Modelo Entidade-Relacionamento (MER) e o Diagrama Entidade-Relacionamento (DER), além da documentação inicial sobre regras de negócio e segurança da informação.  

O objetivo principal dessa etapa foi consolidar toda a base teórica e técnica para o início da codificação do sistema.

\subsubsection{Fase 2 – Prova de Conceito (POC) (10/2025)}

A Prova de Conceito teve como finalidade validar a comunicação entre o \textit{front-end} e o \textit{back-end}, bem como a integração com o banco de dados MySQL. Nesta etapa, foram realizados testes iniciais de conexão, autenticação de usuários e manipulação de dados para garantir o funcionamento do ambiente de desenvolvimento.  

Com isso, foi possível comprovar a viabilidade técnica da solução proposta.

\subsubsection{Fase 3 – Produto Mínimo Viável (MVP) (10/2025 – 11/2025)}

Nesta fase, foi implementada uma versão funcional do sistema, contendo as principais operações: controle de estoque, registro de vendas, autenticação de usuários. 

O MVP serviu como base para avaliações práticas junto ao contexto da loja, possibilitando ajustes antes da entrega final.

\subsubsection{Fase 4 – Entrega Final (Data à definir)}

A entrega final compreende a versão completa e refinada do sistema, com todas as funcionalidades implementadas, correções aplicadas e testes de desempenho concluídos. Essa etapa também incluiu a revisão dos requisitos de segurança e conformidade com a LGPD, além da documentação técnica e do manual de uso.  

